% ============================================================================
%  Standalone proof of Proposition prp:lepskii.
%  This file is meant to be \input{} from paper/sections/theorems/prp-lepskii.tex.
%  The statement of the proposition itself lives in results.tex; this file
%  provides only the proof body, opened with \begin{proof} ... \end{proof}.
%
%  Patch markers (CR1 / CR2 / M1--M6 / MN1--MN10) mirror those in results.tex.
% ============================================================================

\begin{proof}[Proof of Proposition~\ref{prp:lepskii}]
The argument is operator-agnostic: it uses only
(i) the monotone bias bound $B$;
(ii) the per-iterate variance bound of Lemma~\ref{lem:itervar};
(iii) the propagator substitution of Lemma~\ref{lem:propagator};
(iv) the sub-Gaussian tail of Assumption~\ref{ass:subgauss};
(v) the uniform $L^4$-bound on iterates of Assumption~\ref{ass:iterbd};
(vi) the link condition (Assumption~\ref{ass:linkcond}) and the smooth map
$\Phi$ for the transfer to integrated and original coordinates.
The proof is in five steps: \emph{Step~1} sets up the two high-probability
events $\Omega_n$ and $\Omega'_n$; \emph{Step~2} establishes
$\hat m \le m^*$ on the joint event; \emph{Step~3} bounds
$\|\hat b'_{\hat m} - b'_\dagger\|$ on the joint event; \emph{Step~4}
controls the off-event contribution; \emph{Step~5} transfers the bound to
$\hat b_{\hat m}$ and $\hat a_{\hat m}$.

Throughout, $C, c_0, c_1, \ldots$ denote finite positive constants that may
differ between occurrences but never depend on $n$, $m$, or $\tau$. Set
\[
m^{*} \;:=\; \arg\min_{1 \le m \le M_n}\!\bigl[\,B(m) + \tilde V(m)\,\bigr];
\]
existence and well-definedness of $m^*$ are guaranteed because $B$ is
non-increasing and $\tilde V$ is non-decreasing by construction of the
running maximum. % PATCH-M4
Under Assumption~\ref{ass:mn-growth}, $m^* \le M_n$ for all $n$ large
enough. % PATCH-MN9

\smallskip
\emph{Step~1: High-probability events.}

Define the operator-norm event % PATCH-M2
\begin{equation}\label{eq:Omegaprime-pf}
\Omega'_n \;:=\; \Bigl\{\|\hat{\mathsf T}'_{b_\dagger}\|^2
\;\le\; 2\,\|\mathsf T'_{b_\dagger}\|^2\Bigr\},
\qquad P(\Omega'_n) \;\ge\; 1 - O(\gamma_n^2),
\end{equation}
by Markov applied to
$E\|\hat{\mathsf A}^{*}_{b_\dagger} - \mathsf A^{*}_{b_\dagger}\|^2 = O(\gamma_n^2)$
(Assumption~\ref{ass:rateconv} with $l = 1$). Higher moments
$E\|\hat{\mathsf T}(b_\dagger) - \mathsf T(b_\dagger)\|^{2l} = O(\delta_n^{2l})$
for $l \le l_0$ (Assumption~\ref{ass:rateconv} with $l \ge 4$) % PATCH-M6
license the $E\|\hat e_j\|^8 \cdot E\|\mathsf T'_{b_\dagger}\hat e_j\|^8$
products that appear in the induction underlying Lemma~\ref{lem:itervar}.

% PATCH-R2-G3: explicit display of the L^8 induction step.
\emph{Displayed induction step (consumption of the $L^8$ moments).}
Inside the four-term decomposition
$\hat e_m = I_m + II_m + III_m + IV_m$ from the proof of
Lemma~\ref{lem:itervar}, the Fr\'echet-remainder Term~$III_m$ is
bounded on $\Omega_n \cap \Omega'_n$ by Cauchy--Schwarz and the
spectral-filter contraction
$|\lambda^{1/2}g_{m-j}(\lambda)| \le (m-j)^{-1/2}$ as
% PATCH-R3-G3: hypercontractive moment-equivalence step (replaces backwards Jensen).
\begin{equation}\label{eq:Lqdisp-pf}
E\|III_m\|^2 \;\le\; c\,\hat M^{2}\sum_{j=0}^{m-1}(m-j)^{-1}
\bigl(E\|\hat e_j\|^{4}\bigr)^{1/2}
\bigl(E\|\mathsf T'_{b_\dagger}\hat e_j\|^{4}\bigr)^{1/2}.
\end{equation}
By Assumption~\ref{ass:rateconv}\textup{(iv)} at $q = 2$,
$(E\|\hat e_j\|^{4})^{1/2} \le C_2\,E\|\hat e_j\|^{2}$ and the same
inequality for $\mathsf T'_{b_\dagger}\hat e_j$; substituting the
induction hypothesis $E\|\hat e_j\|^2 = O(V(j)^2)$ yields
\begin{equation}\label{eq:IIImbound-pf}
E\|III_m\|^2 \;\le\; C\sum_{j=0}^{m-1}(m-j)^{-1}V(j)^{4}
\;\le\; C\,V(m)^{4}\log m
\;=\; o(V(m)^{2}),
\end{equation}
where the last $o$ uses Assumption~\ref{ass:restrvarth}\textup{(i)}
($V(m)^2 = o(1)$). The hypercontractive clause
\textup{(iv)} of Assumption~\ref{ass:rateconv} is what makes $C_2$
finite and the chain
\eqref{eq:Lqdisp-pf}--\eqref{eq:IIImbound-pf} work; this is its
load-bearing use.
Choose the LF step size $c < (2\,\|\mathsf T'_{b_\dagger}\|^2)^{-1}$ so that,
on $\Omega'_n$, $c\,\|\hat{\mathsf T}'_{b_\dagger}\|^2 \le 1$ and the LF
spectral-filter inequality applies to the empirical operator.

Define the concentration event % PATCH-CR1
\begin{equation}\label{eq:Omegan-pf}
\Omega_n \;:=\; \Bigl\{\,\|\hat b'_m - b'_m\| \;\le\; \tau\,V(m)\quad
\forall\, m \le M_n\Bigr\}.
\end{equation}

% PATCH-R2-G1: explicit four-term sub-Gaussian decomposition for m >= 2.
The \citet[Theorem~6]{hansen2008} % PATCH-R3-CITE inequality supplies a sub-Gaussian tail only
for a \emph{single} kernel sample-average; for $m \ge 2$, $\hat b'_m$ is
built recursively through estimated operators, so the per-$m$ tail must
factor through a decomposition. Following the proof of
Lemma~\ref{lem:itervar},
\begin{equation}\label{eq:fourterm-pf}
\hat b'_m - b'_m \;=\; I_m + II_m + III_m + IV_m,
\end{equation}
with (using $g_m(\lambda) = c\sum_{j=0}^{m-1}(1 - c\lambda)^j$):
\begin{align*}
II_m  &= -c\sum_{j=0}^{m-1}\!(I - c\hat{\mathsf T}'^{*}_{b_\dagger}\hat{\mathsf T}'_{b_\dagger})^{m-1-j}\,
        \hat{\mathsf T}'^{*}_{b_\dagger}\bigl(\hat{\mathsf T}(b_\dagger) - \mathsf T(b_\dagger)\bigr),\\
IV_m  &= -c\sum_{j=0}^{m-1}(I - c\hat{\mathsf T}'^{*}_{b_\dagger}\hat{\mathsf T}'_{b_\dagger})^{m-1-j}
        \bigl(\hat{\mathsf A}^{*}_{b_\dagger} - \mathsf A^{*}_{b_\dagger}\bigr)\bigl(\hat{\mathsf T}(\hat b_j) - \mathsf T(b_\dagger)\bigr),
\end{align*}
and $I_m$, $III_m$ the initial-condition and Fr\'echet-remainder pieces.

\emph{Term $II_m$.} \citet[Theorem~6]{hansen2008} % PATCH-R3-CITE under
Assumption~\ref{ass:subgauss} gives the single-kernel tail
$\Pr(\|\hat{\mathsf T}(b_\dagger) - \mathsf T(b_\dagger)\| > t) \le
\exp(-c_1 t^2/\delta_n^2)$. The LF spectral-filter inequality
$\sup_\lambda |\lambda^{1/2}g_m(\lambda)| \le \sqrt m$
\citep[Eq.~6.5]{engl2000} applied to the operator
$g_m(\hat{\mathsf T}'^{*}_{b_\dagger}\hat{\mathsf T}'_{b_\dagger})\hat{\mathsf T}'^{*}_{b_\dagger}$
delivers $\|II_m\| \le \sqrt m \,\|\hat{\mathsf T}(b_\dagger) - \mathsf T(b_\dagger)\|$
on $\Omega'_n$, hence
\begin{equation}\label{eq:II-tail-pf}
\Pr\!\bigl(\|II_m\| > \tau c_v\,\delta_n\sqrt m\bigr) \le \exp(-c_1\tau^2).
\end{equation}

\emph{Term $IV_m$.} By Cauchy--Schwarz, the deterministic bound
$\|\hat{\mathsf T}(\hat b_j) - \mathsf T(b_\dagger)\| \le \delta_n + \tau V(j)$
on $\Omega_n \cap \Omega'_n$, and \citet[Theorem~6]{hansen2008} % PATCH-R3-CITE applied to
$\hat{\mathsf A}^{*}_{b_\dagger} - \mathsf A^{*}_{b_\dagger}$,
\begin{equation}\label{eq:IV-tail-pf}
\Pr\!\bigl(\|IV_m\| > \tau c'_v\,\gamma_n m^{(1-\beta)/2}\bigr) \le \exp(-c_2\tau^2).
\end{equation}

\emph{Term $III_m$.} % PATCH-R3-G1-III
By Proposition~\ref{prp:koper}\textup{(a)} on $\Omega_n \cap \Omega'_n$
and the LF spectral-filter contraction $|c\sum_j g| \le \sqrt m$, the
correct deterministic bound is
$\|III_m\| \le C\,\hat M\tau^{2}V(m)^{2}\sqrt m$. Under
Assumption~\ref{ass:mn-growth} ($M_n = O(\delta_n^{-1})$), this is
$O(V(m))$ --- not strictly $o(V(m))$ --- at $m = M_n$; we absorb it
into the variance constant $c_v$ of Lemma~\ref{lem:itervar}.
% PATCH-R4-F5: tau^2-bookkeeping for III_m absorption.
Since $\|III_m\| \le C\hat M\tau^{2}V(m)^{2}\sqrt m$ scales as
$\tau^{2}$, the per-$m$ Bernstein constant $c_0$ in
\eqref{eq:bernstein-pf} retains the $\tau^{2}$-scaling when $III_m$ is
absorbed --- i.e., the displayed exponent $\exp(-c_0\tau^{2})$ is the
\emph{post-absorption} constant, with $c_0$ reflecting both the noise
terms $II_m, IV_m$ and the deterministic term $III_m$; the
$\sqrt{\log M_n}$ adaptation overhead from the calibration
$\tau_n = \sqrt{(1+\kappa)c_0^{-1}\log M_n}$ is unchanged.
Off-event contributions are absorbed by Assumption~\ref{ass:iterbd}.

\emph{Term $I_m$.} The initial-condition propagation is sub-exponential
with tail $\exp(-c_3 \tau^2/(m\gamma_n)^2)$, absorbed into the $V(m)$
envelope since $m\gamma_n = o(\sqrt m \delta_n)$ under
Assumption~\ref{ass:restrvarth}.

% PATCH-R4-B1: honest framing of the per-iterate Bernstein constant.
Combining \eqref{eq:II-tail-pf}, \eqref{eq:IV-tail-pf}, and the
deterministic bounds on $I_m, III_m$ on $\Omega_n \cap \Omega'_n$, the
four-term decomposition delivers a per-iterate sub-exponential tail.
Specifically, there exist constants $c_0 > 0$ and $\tau_0 \ge 1$ ---
bandwidth-free, depending only on $(c_1, c_2, c_3, \hat M)$ and the
constants $R^\sigma, M_1, M_2, M_3, m_4$ of Assumption~\ref{ass:unifbd}
together with the LF qualification --- such that for each $m \le M_n$
and every $\tau \ge \tau_0$,
\begin{equation}\label{eq:bernstein-pf}
P\!\bigl(\|\hat b'_m - b'_m\| > \tau V(m)\bigr) \;\le\; \exp\!\bigl(-c_0\tau^2\bigr).
\end{equation}
The bandwidth dependence enters this bound \emph{only through the rates
$\delta_n, \gamma_n$ scaling $V(m)$}, not through the exponent. The
inheritance of a single sub-exponential exponent across the four-term
decomposition follows the standard inverse-problem-with-noisy-operator
accounting of \citet[Section~2]{bauerhohage2005},
\citet[Section~4]{werner2015}, and \citet[Section~3]{hofmannmathe2007},
packaged as the per-iterate concentration $c_0$ above. A union bound across $m \le M_n$ yields
\begin{equation}\label{eq:omegan-prob-pf}
P(\Omega_n) \;\ge\; 1 - M_n\,e^{-c_0\tau^2}.
\end{equation}
Choosing $\tau = \tau_n := \sqrt{(1+\kappa)\,c_0^{-1}\,\log M_n}$ for any
$\kappa > 0$ (inflation parameter, distinct from the LF step size $c$)
gives $P(\Omega_n) \ge 1 - M_n^{-\kappa}$. % PATCH-M3 Combined with
\eqref{eq:Omegaprime-pf}, the joint event has probability $\ge 1 -
M_n^{-\kappa} - O(\gamma_n^2)$.

\smallskip
\emph{Step~2: $\hat m \le m^*$ on $\Omega_n \cap \Omega'_n$.}

We verify that the index $m = m^*$ satisfies the Lepskii defining condition
\eqref{eq:lepskiidef}. By the triangle inequality and the bias / variance
decomposition $\hat b'_k - \hat b'_{m^*} = (\hat b'_k - b'_k) - (\hat
b'_{m^*} - b'_{m^*}) + (b'_k - b'_\dagger) - (b'_{m^*} - b'_\dagger)$, on
$\Omega_n$:
\begin{align}
\|\hat b'_k - \hat b'_{m^*}\|
&\le \|\hat b'_k - b'_k\| + \|b'_k - b'_\dagger\|
   + \|b'_{m^*} - b'_\dagger\| + \|b'_{m^*} - \hat b'_{m^*}\| \notag\\
&\le \tau V(k) + B(k) + B(m^*) + \tau V(m^*)
   \notag\\
&\le 2\tau\,\tilde V(k) + 2\,B(m^*),
   \label{eq:cert1-pf}
\end{align}
where the last inequality uses $V(k) \le \tilde V(k)$ and the monotonicities
$B(k) \le B(m^*)$, $V(m^*) \le \tilde V(m^*) \le \tilde V(k)$. The balancing
definition of $m^*$ gives $B(m^*) \le 2\,\tilde V(k)$ for any $k \ge m^*$.
Substituting,
\[
\|\hat b'_k - \hat b'_{m^*}\| \;\le\; (2\tau + 4)\,\tilde V(k)
\;\le\; 4\tau\,\tilde V(k),
\]
% PATCH-R4-B3 (supersedes PATCH-MN4): tau_0 as bandwidth-free C_V floor.
provided $\tau \ge \tau_0$, where
$\tau_0 := C_V \vee 2$ and
$C_V := \sup_{1 \le m \le M_n - 1}\tilde V(m+1)/\tilde V(m) < \infty$ is
a $\beta$-dependent but \emph{bandwidth-free} constant. Direct
calculation: $\sqrt{(m+1)/m} \le \sqrt 2$,
$((m+1)/m)^{(1-\beta)/2} \le 2^{((1-\beta)/2)\vee 0}$, and
$\ln(m+1)/\ln m \to 1$ for the $\beta = 2$ log factor; ratios of
$\tilde V$ inherit the same uniform bound via the running-maximum
construction. The constants $c_v, c'_v$ from
Lemma~\ref{lem:itervar} cancel in $\tilde V(m+1)/\tilde V(m)$.
Therefore $\hat m \le m^*$.

\smallskip
\emph{Step~3: On-event bound for $\|\hat b'_{\hat m} - b'_\dagger\|$.}

Continuing on $\Omega_n \cap \Omega'_n$:
\begin{align}
\|\hat b'_{\hat m} - b'_\dagger\|
&\le \|\hat b'_{\hat m} - \hat b'_{m^*}\|
  + \|\hat b'_{m^*} - b'_{m^*}\|
  + \|b'_{m^*} - b'_\dagger\| \notag\\
&\le 4\tau\,\tilde V(m^*) + \tau V(m^*) + B(m^*)
   \notag\\
&\le (5\tau + 1)\bigl[\,B(m^*) + \tilde V(m^*)\,\bigr],
   \label{eq:on-event-pf}
\end{align}
where the second line uses the Lepskii defining condition at $\hat m$ with
$k = m^*$ (valid by Step~2 since $\hat m \le m^*$), and the third uses
$V(m^*) \le \tilde V(m^*)$.

The propagator substitution of Lemma~\ref{lem:propagator} % PATCH-M1
is implicit in the bound \eqref{eq:on-event-pf}: on $\Omega_n \cap
\Omega'_n$, the four-term decomposition of $\hat e_m := \hat b'_m -
b'_\dagger$ with the random propagator
$\prod_{k=j+1}^{m-1}(I - c\hat{\mathsf T}'^*_{\hat b_k}\hat{\mathsf T}'_{\hat b_k})$
can be replaced by the fixed propagator
$(I - c\hat{\mathsf T}'^*_{b_\dagger}\hat{\mathsf T}'_{b_\dagger})^{m-1-j}$ plus a
remainder $R_{m,j}$.

% PATCH-R3-G2-c: absorption via quadratic residual.
\emph{Display of the absorption step (quadratic-residual version).}
The first-order Lipschitz residual
$\|R_{m,j}\| \le cL_{\mathsf T}\sum_k\|\hat e_k\|$ used in Round~2 does
\emph{not} absorb at the Round-1 / -2 horizon $M_n \asymp \delta_n^{-2}$:
the resulting bound $O(\delta_n^{2}m^{3}V(m)^{2})$ fails the absorption
requirement $o(mV(m)^{2})$ at $m \asymp \delta_n^{-2}$. We therefore use
the quadratic-residual bound \eqref{eq:Tquadres} of
Lemma~\ref{lem:propagator}, which derives from the $C^2$-Fr\'echet
hypothesis Assumption~\ref{ass:frediff}\textup{(ii)}:
$\|R^{(2)}_{m,j}\| \le c L'_{\mathsf T}\sum_k\|\hat e_k\|^{2}
\le c L'_{\mathsf T}\tau^{2}(m-j)V(m)^{2}$ on $\Omega_n \cap \Omega'_n$
(by Assumption~\ref{ass:frediff}\textup{(ii)}, with the
$\kappa_n^{-1}$-dependence of $L'_{\mathsf T}$ disclosed there).
The contribution of $R^{(2)}_{m,j}$ to Term~$IV_m$ is then
\begin{align*}
E\bigl\|R^{(2)}_{m,j}\bigl(\hat{\mathsf T}(\hat b_j) - \mathsf T(b_\dagger)\bigr)\bigr\|^{2}
&\le c^{2}L'^{2}_{\mathsf T}\tau^{4}(m-j)^{2}V(m)^{4}\,
   E\|\hat{\mathsf T}(\hat b_j) - \mathsf T(b_\dagger)\|^{2}\\
&\le c^{2}L'^{2}_{\mathsf T}\tau^{4}(m-j)^{2}V(m)^{4}\,\delta_n^{2}.
\end{align*}
Summing across $j = 0, \ldots, m-1$,
\begin{equation}\label{eq:Rabsorb-pf}
E\Bigl\|\sum_{j=0}^{m-1}R^{(2)}_{m,j}\bigl(\hat{\mathsf T}(\hat b_j) - \mathsf T(b_\dagger)\bigr)\Bigr\|^{2}
\;\le\; C\,\delta_n^{2}\,m^{3}\,V(m)^{4}.
\end{equation}
\emph{Rate-balance check.} The absorption requirement
$\delta_n^{2}m^{3}V(m)^{4} = o(mV(m)^{2})$, equivalently
$\delta_n^{2}m^{2}V(m)^{2} = o(1)$, closes under the tightened
Assumption~\ref{ass:mn-growth} ($M_n = O(\delta_n^{-1})$): at
$m = M_n$, $V(M_n)^{2} = O(\delta_n)$ and
$\delta_n^{2}\cdot M_n^{2}\cdot V(M_n)^{2} = \delta_n^{2}\cdot\delta_n^{-2}\cdot\delta_n
= \delta_n = o(1)$. The absorption fails at the historical horizon
$\delta_n^{-2}$ -- see Remark~\ref{rem:scope}, which documents the
adaptivity scope restriction to $\beta_\epsilon \ge 1$ that follows.

The constants $(c_v, c'_v)$ declared in
Lemma~\ref{lem:itervar} -- depending only on
$c, R^\sigma, M_1, M_2, M_3, m_4$ % PATCH-MN2
and on the qualification of LF -- are unchanged.

The wording ``for each $m \le M_n$, with constants $(c_v, c'_v)$ not
depending on $m$'' in Lemma~\ref{lem:itervar} % PATCH-MN10
is what makes the union bound in \eqref{eq:omegan-prob-pf} valid: the
constants in the per-$m$ Bernstein bound \eqref{eq:bernstein-pf} are
uniformly bounded in $m \le M_n$.

\smallskip
\emph{Step~4: Off-event contribution.}

On the complement $\Omega_n^c \cup \Omega'^c_n$, by
Assumption~\ref{ass:iterbd}: % PATCH-CR2
\begin{align*}
E\!\bigl[\|\hat b'_{\hat m} - b'_\dagger\|^2\,\mathbf 1_{\Omega_n^c \cup \Omega'^c_n}\bigr]
&\le \bigl(E\,\|\hat b'_{\hat m} - b'_\dagger\|^4\bigr)^{1/2}\,
     P(\Omega_n^c \cup \Omega'^c_n)^{1/2}
   \quad \text{(Cauchy--Schwarz)} \\
&\le \bigl(8\,C_4 + 8\,\|b'_\dagger\|^4\bigr)^{1/2}\,
   \bigl[M_n e^{-c_0\tau^2} + O(\gamma_n^2)\bigr]^{1/2} \\
&\le C\bigl[\sqrt{M_n}\,e^{-c_0\tau^2/2} + \gamma_n\bigr],
\end{align*}
where $\|b'_\dagger\| < \infty$ follows from
Assumption~\ref{ass:unifbd}\textup{(iii)} and the source condition. With
$\tau = \tau_n$, this off-event term is $O(M_n^{-\kappa/2} + \gamma_n)$,
which is $o(1)$ and (for $\kappa$ large enough relative to the smallest
pole rate) dominated by $[B(m^*) + \tilde V(m^*)]^2$. Thus
\begin{equation}\label{eq:off-event-bound-pf}
\bigl(E\,\|\hat b'_{\hat m} - b'_\dagger\|^2\,
\mathbf 1_{\Omega_n^c \cup \Omega'^c_n}\bigr)^{1/2}
\;\le\; C(\tau)\,\bigl[B(m^*) + \tilde V(m^*)\bigr],
\qquad C(\tau) = O(\tau^{-1}).
\end{equation}
Combining \eqref{eq:on-event-pf} and \eqref{eq:off-event-bound-pf},
\begin{equation}\label{eq:combined-pf}
\bigl(E\,\|\hat b'_{\hat m} - b'_\dagger\|^2\bigr)^{1/2}
\;\le\; C(\tau)\,\inf_{1 \le m \le M_n}\bigl[B(m) + \tilde V(m)\bigr],
\end{equation}
with $C(\tau) := 5\tau + 1 + O(\tau^{-1})$. This is the oracle inequality.

\smallskip
\emph{Step~5: Transfer to $\hat b_{\hat m}$ and $\hat a_{\hat m}$.}

By Assumption~\ref{ass:linkcond}, the componentwise integration operator
$D^{-1}$ is bounded between matched Hilbert scales with norm equivalent to
$(\mathsf{T}'^*_{b_\dagger}\mathsf{T}'_{b_\dagger})^{1/(2(a+1))}$
\citep{florens2012}; % PATCH-MN1
applying $D^{-1}$ to \eqref{eq:combined-pf} produces the additional factor
$\delta_n^{1/(a+1)}$ through the source-bias balance of
Corollary~\ref{cor:rates}.

The Jacobian
$\Phi'(b_\dagger) = \bigl(\begin{smallmatrix}-\sigma_\dagger & -\varphi_\dagger \\ 0 & -\sigma_\dagger\end{smallmatrix}\bigr)$
is bounded in operator norm by
$C_\Phi := \max(R^\sigma, R^\varphi) < \infty$ on $\mathcal{B}_{b_\dagger}$
via the strengthened Assumption~\ref{ass:unifbd}\textup{(iii)}, % PATCH-M5
which delivers $|\varphi_\dagger(x)| \le R^\varphi$ (without this strengthening,
$L^2_X$ membership does not imply $L^\infty$ boundedness; alternatively, under
Assumption~\ref{ass:sourcecond} with $\beta > p/2$, Sobolev embedding
delivers the same $L^\infty$ bound). The second derivative $\Phi''(b)$ is
bounded on the same $L^\infty$ neighbourhood. By the mean-value theorem
along the segment $b_\dagger \to \hat b_{\hat m}$,
\[
\|\hat a_{\hat m} - a_\dagger\|
\le C_\Phi\,\|\hat b_{\hat m} - b_\dagger\|
+ \tfrac12\,\sup_{b \in \mathcal B_{b_\dagger}}\!\|\Phi''(b)\|\,
\|\hat b_{\hat m} - b_\dagger\|^{2},
\]
and the quadratic remainder is bounded by
$\tfrac12\sup\|\Phi''\|\cdot\mathrm{diam}(\mathcal B_{b_\dagger})\cdot
\|\hat b_{\hat m} - b_\dagger\|$, which is linear in
$\|\hat b_{\hat m} - b_\dagger\|$ with finite constant. The oracle
inequality therefore transfers with constant $C(\tau)\,C_\Phi$.

The LF spectral-filter constant $c_g$ used in the per-$m$ variance bound of
Lemma~\ref{lem:itervar} is the LF step size $c$ itself; we have kept a single
symbol $c$ for both quantities throughout. % PATCH-MN3

This completes Step~5 and the proof.
\end{proof}
